\subsection{Introduzione}
\label{subsec:introduzione}

Il sistema ArtificialQI è basato sull'architettura client-server.
In questa architettura il carico di lavoro viene suddiviso tra due componenti:
\begin{enumerate}
    \item \textbf{client}: interfaccia con cui gli utenti interagiscono con il sistema per richiedere risorse e/o servizi;
    \item \textbf{server}: applicazione che riceve ed elabora le richieste eseguite dai client e fornisce le risposte corrette. 
\end{enumerate}
Il gruppo ha quindi deciso di sviluppare due componenti separati:
\begin{enumerate}
    \item \textbf{front-end}: componente client del sistema, è un interfaccia grafica sviluppata utilizzando il framework Angular;  
    \item \textbf{back-end}: componente server del sistema che utilizzando il framework Flask implementa la logica di business e la logica di persistenza.
\end{enumerate}
La comunicazione tra front-end e back-end avviene tramite una API.
L'API verrà implementata seguendo i principi indicati nello stile architetturale REST (Representational State Transfer), utilizzando il protocollo di comunicazione HTTP e il formato di rappresentazione delle risorse JSON.
Una caratteristica delle REST API è quella di essere stateless, il che significa che la componente server non deve memorizzare le interazioni passate per poter rispondere a una richiesta del client. 
Questo permette di ridurre il carico sul server e migliora la scalabilità.
La persistenza delle risorse viene implementata utilizzando un database relazionale gestito utilizzando PostgreSQL come DBMS.

\subsubsection{Integrazione dei componenti}
Il gruppo ha deciso di utilizzare un container Docker per ognuno dei due componenti sopra citati.
Per gestire e permettere la comunicazione tra i due componenti il gruppo ha deciso di utilizzare Docker Compose.  
